\subsection{The sparse tools display different normalized structures across models}
\label{sec:sparse-structure}

The preceding same-object intervention shows that the fitted sparse reconstructions do not reliably recover target-answer scores. Their feature patterns may nevertheless describe how each fitted tool responds to the image corruption. We therefore ask how the clean--corrupted change is distributed within each sparse representation. Because feature identities are not shared across models, we compare normalized summaries rather than raw feature numbers, and we do not use these summaries to infer a faithful path.

\begin{samepage}
The clean--masked support overlap is the fraction of active feature identities shared by the two image conditions. The feature fraction for 80\% of change is the proportion of the tool's full feature set needed to account for 80\% of the squared sparse change; a smaller value means that the change is more concentrated. The positive change-energy fraction measures how much of that squared change comes from features that increase rather than decrease. Finally, the log sparse/full norm ratio compares the size of the sparse reconstruction with the complete internal change within the same model and position group. Table~\ref{tab:sparse-structure-summary} reports these four summaries across models and position groups.
\end{samepage}

\begin{samepage}
\begin{center}
\begin{minipage}{\linewidth}
\captionsetup{type=table,hypcap=false}
\caption{\textbf{The structures seen by the sparse tools differ across models.} Values are model means on the same 240 images. ``Pairs'' counts how many of the three paired model comparisons remain different after the joint correction of 72 tests. These summaries describe each fitted tool and do not align feature identities across models.}
\label{tab:sparse-structure-summary}
\centering
\small
\setlength{\tabcolsep}{2.2pt}
\begin{tabular*}{\linewidth}{@{\extracolsep{\fill}}llrrrr@{}}
\toprule
Summary & Positions & Gemma & Qwen & LLaVA & Pairs \\
\midrule
Clean--masked support overlap & Visual & 0.351 & 0.406 & 0.407 & 2/3 \\
Clean--masked support overlap & Answer & 0.677 & 0.690 & 0.684 & 0/3 \\
Feature fraction for 80\% of change & Visual & 0.0475 & 0.0660 & 0.0345 & 3/3 \\
Feature fraction for 80\% of change & Answer & 0.0142 & 0.0099 & 0.0083 & 3/3 \\
Positive change-energy fraction & Visual & 0.512 & 0.525 & 0.510 & 1/3 \\
Positive change-energy fraction & Answer & 0.491 & 0.495 & 0.519 & 2/3 \\
Log sparse/full norm ratio & Visual & -0.159 & -0.247 & -0.151 & 2/3 \\
Log sparse/full norm ratio & Answer & 0.079 & 0.125 & -0.001 & 3/3 \\
\bottomrule
\end{tabular*}
\end{minipage}
\end{center}
\end{samepage}


Sixteen of the 24 paired model comparisons remain different after all 72 tests are corrected together. The clearest pattern is concentration. At visual positions, the feature fractions needed for 80\% of the change are $.0660$ for Qwen, $.0475$ for Gemma, and $.0345$ for LLaVA; at answer positions, they are $.0142$ for Gemma, $.0099$ for Qwen, and $.0083$ for LLaVA. All six pairwise comparisons for this summary remain different after correction. The other summaries are less uniform: clean--masked support overlap differs for two visual-position pairs but no answer-position pair, while positive change energy differs for one visual-position pair and two answer-position pairs.

These results establish model-conditioned structural differences in what the fitted tools display. They do not establish a shared feature identity, and they do not show that the base models themselves use the same or different complete sparse mechanisms. The earlier intervention found no reliable positive target-score change from these reconstructions despite their different normalized structures. We next remove the fitted tools from the intervention and test whether selected channels belonging to the base models account for the score change.

\FloatBarrier
